\documentclass[11pt,a4paper]{article}
\usepackage[utf8]{inputenc}
\usepackage{amsmath,amssymb}
\usepackage{geometry}
\geometry{margin=1in}

\title{\textbf{Computational Chemistry \& Molecular Modeling}}
\author{C S Tarun Kumar}
\date{November 20, 2026}

\begin{document}
\maketitle

\section{Overview}
Computational chemistry utilizes computer simulation to solve chemical problems, calculate molecular structures, and predict physical properties.

\section{Key Approaches}
\subsection{Molecular Mechanics (MM)}
\begin{itemize}
    \item Treats atoms as balls and bonds as springs (force field methods).
    \item Fast and computationally inexpensive; ideal for large biomolecules.
    \item Does not account for electronic transitions or bond formation/breaking.
\end{itemize}

\subsection{Electronic Structure Methods (Quantum Mechanics)}
\begin{itemize}
    \item \textbf{Ab Initio:} Solves the Schrödinger equation from first principles (Hartree-Fock, Post-HF methods like MP2, CCSD).
    \item \textbf{Density Functional Theory (DFT):} Calculates electron density rather than multi-electron wavefunctions (e.g., B3LYP functional). Excellent balance of accuracy and computational cost.
    \item \textbf{Semi-Empirical Methods:} Uses empirical parameters to approximate complex integrals (AM1, PM3).
\end{itemize}

\section{Practical Applications}
\begin{itemize}
    \item Drug discovery and protein-ligand docking.
    \item Transition state analysis and reaction kinetics.
    \item Spectroscopic property prediction (IR, NMR, UV-Vis).
\end{itemize}

\end{document}
